% BANKS-SOURCE: pdf=54 print=44 kind=chapter id=chapter05
\BanksChapter{5}{Spin-$\frac12$ particles and Fermi statistics}

% BANKS-SOURCE: pdf=54 print=44 kind=prose
For spin-$\frac12$ particles, we will introduce a change of pace and start directly from field
theory. We have to ask ourselves what kind of fields could create spin-$\frac12$ particles. This
leads to the question of what kinds of fields there are, which is answered by saying that
fields at a fixed point form finite-dimensional representations of the Lorentz group.
This guarantees that the only infinity in the number of states of a single particle comes
from the different values of momentum it can carry.

We’re lucky to live in four dimensions, where the analysis of finite-dimensional
representations is particularly easy. In particular, the kind reader will verify that the
combinations of Lorentz generators\footnote[1]{We are picking a particular rotation subgroup of $\mathrm{SO}(1,3)$ in order to exhibit the invariant isomorphism between groups.}
\begin{equation*}
  \epsilon_{ijk}J_{ij}\pm\ii J_{0k}
\end{equation*}
where each independent spatial pair is counted once, form two commuting copies of the algebra of $\mathrm{SU}(2)$. We can use our complete knowl-
edge of the representations of $\mathrm{SU}(2)$ to list all of the finite-dimensional representations
of $\mathrm{SO}(1,3)$. The finite-dimensional representations of $\mathrm{SU}(2)$ are all equivalent to uni-
tary representations, so we can always use a basis in which the two $\mathrm{SU}(2)$ generators
are Hermitian. Note that, using this construction, the rotation generators will be Her-
mitian and the boost generators anti-Hermitian. This was to be expected. We cannot
have a finite-dimensional unitary representation of the non-compact Lorentz group.\footnote[2]{It is important to remember that this is the representation of the Lorentz group on field labels. The representation on the Hilbert space is infinite-dimensional and unitary.}
Thus, the general Lorentz-covariant field carries two integer labels $N\equiv[n_L,n_R]$ corre-
sponding to the dimensions of the representations of the two $\mathrm{SU}(2)$ groups. Recall that
even dimensions correspond to half-integer and odd to integer spin. Note that the two
kinds of integer are interchanged by space reflection. Thus, fields of the form $[1,n_R]$
and $[n_L,1]$ have a handedness, and are called right and left chiral fields.

\BanksImplicitHook{I-CH05-001}

% BANKS-SOURCE: pdf=54 print=44 kind=prose
The rotation generators are the sum of the left and right $\mathrm{SU}(2)$ generators, so we can
get spin $\frac12$ with either $[1,2]$ or $[2,1]$. These are called right- and left-handed Weyl spinors,
respectively. In fact the $[1,2]$ representation is complex,\footnote[3]{$\mathrm{SO}(1,3)$ is locally isomorphic to $\mathrm{SL}(2,\mathbb C)$ and the $[1,2]$ is the fundamental representation of the latter.} and its complex conjugate is

% BANKS-SOURCE: pdf=55 print=45 kind=prose
the $[2,1]$. To see this note that the rotation and boost generators in the $[1,2]$ are $J_i=
\sigma_i/2$, $J_{0i}=-\ii\sigma_i/2$. The representation of rotations is pseudo-real, $\sigma_i^*=-\sigma_2\sigma_i\sigma_2$, but
the boost operator has an extra sign change under conjugation, so that we get the $[2,1]$
representation. We can describe a spin-$\frac12$ particle either in terms of a left-handed Weyl
fermion field $\nu$, in the $[2,1]$ representation, and its complex conjugate, or equivalently,
in terms of the four-component field\footnote[4]{In this section we will use the symbol $*$ to denote Hermitian conjugation of the operators in each component of a spinor field. This is to distinguish it from $\dagger$, which could be taken to imply turning row indices into column indices as well as Hermitian conjugation. Later on, most of our equations will be written in terms of complex Dirac spinors, and we will revert to the usual symbol for Hermitian conjugation.}
\begin{equation*}
  \psi=\begin{pmatrix}\nu\\ \ii\sigma_2\nu^*\end{pmatrix},
\end{equation*}
which is a Dirac spinor field, satisfying the Majorana condition
\begin{equation*}
  \psi^*=\begin{pmatrix}0&-\ii\sigma_2\\ \ii\sigma_2&0\end{pmatrix}\psi.
\end{equation*}
The equivalence between Weyl and Majorana descriptions of a neutral spin-$\frac12$ particle
was not always recognized, and there is a lot of confusion in the literature, as late as the
1970s. A general Dirac spinor has the form $\psi_1+\ii\psi_2$, with $\psi_i$ satisfying the Majorana
condition.

Following van der Waerden, we describe left-handed Weyl fields as two-component
fields $\nu_a$, whereas right-handed fields are written as $\psi_{\dot a}$. Each of these representations
has a Lorentz-invariant product
\begin{equation*}
  \nu_a^{(1)}\nu_b^{(2)}\epsilon^{ab}\,\psi_{\dot a}^{(1)}\psi_{\dot b}^{(2)}\epsilon^{\dot a\dot b}.
\end{equation*}

% BANKS-SOURCE: pdf=55 print=45 kind=display id=5.0a
We use the Levi-Civita $\epsilon$ symbols\footnote[5]{$\epsilon^{21}=-\epsilon^{12}=1$, and the same for dotted indices.} to raise spinor indices $\nu^a=\epsilon^{ab}\nu_b$ (note that we
contract on the second index). We also introduce $\epsilon_{ab}$ by
\begin{equation*}
  \epsilon_{ab}\epsilon^{bc}=\delta_a^c,
\end{equation*}
and similarly for dotted indices.

The product $[2,1]\otimes[1,2]$ of the two Weyl representations of opposite chirality is
the $[2,2]$ or 4-vector representation. In $D=4$, the matrices
$\sigma^\mu{}_{a\dot b}=(1,\boldsymbol\sigma)_{a\dot b}$ are the Clebsch--
Gordan coefficients relating the tensor product of the two opposite-chirality Weyl
spinor representations to the conventional basis in the 4-vector representation. These
matrices map the right-handed spinor representation into the left-handed spinor. The
corresponding map in the opposite direction is given by $\bar\sigma^\mu{}_{\dot b a}=(1,-\boldsymbol\sigma)_{\dot b a}$. These symbols
allow us to write the Lagrangian

% BANKS-SOURCE: pdf=55 print=45 kind=equation id=5.1
\begin{equation}
  \mathcal{L}_{\mathrm{Weyl}}=-\frac12\ii(\nu^*)^{\dot b}(\bar\sigma^\mu){}_{\dot b a}\partial_\mu\nu^a+\mathrm{h.c.}
  \label{eq:5.1}
\end{equation}

% BANKS-SOURCE: pdf=56 print=46 kind=prose
The two terms are equal, up to a total divergence, and it is conventional to simply drop
both the factor of one half and the Hermitian conjugate term in the Weyl Lagrangian.
If we vary it with respect to $\nu^*$ we get
\begin{equation*}
  -\ii\bar\sigma^\mu\partial_\mu\nu=0,
\end{equation*}
or in momentum space
\begin{equation*}
  p_\mu\bar\sigma^\mu\nu(p)=0.
\end{equation*}
This tells us that positive-energy solutions have negative helicity. The opposite
correlation is valid for right-handed spinors.

The Weyl matrices satisfy
\begin{equation*}
  \sigma_\mu\bar\sigma_\nu=-\eta_{\mu\nu}+\bar\eta^a{}_{\mu\nu}\sigma_a,
  \qquad
  \bar\sigma_\mu\sigma_\nu=-\eta_{\mu\nu}+\eta^a{}_{\mu\nu}\sigma_a.
\end{equation*}
$\sigma$ ($\bar\sigma$) maps the right (left)-handed spinor into the left (right)-handed
one, so the first product maps the left-handed spinor to itself and the second maps the right-handed
spinor to itself. The following self-duality statements are specific to $D=4$, with
$\epsilon^{0123}=+1$. The product of the $[1,2]$ with itself is $[1,1]\oplus[1,3]$, and the second of
these is the anti-self-dual second-rank anti-symmetric tensor $\epsilon_{\mu\nu\alpha\beta}T^{\alpha\beta}=-2T_{\mu\nu}$. $\bar\eta$ is
the Lorentzian continuation of a standard basis for Euclidean anti-self-dual tensors,
introduced by ’t Hooft \cite{banks-ref-25}, which we will study in the last chapter. $\eta$ is the analog for
self-dual tensors. In Lorentzian signature, with $\eta=\operatorname{diag}(-,+,+,+)$,
\begin{equation*}
  \eta^a{}_{\mu\nu}=(\delta_{\nu0}\delta^a_\mu-\delta_{\mu0}\delta^a_\nu)-\ii\epsilon^a{}_{\mu\nu},
  \qquad
  \bar\eta^a{}_{\mu\nu}=(\delta_{\mu0}\delta^a_\nu-\delta_{\nu0}\delta^a_\mu)-\ii\epsilon^a{}_{\mu\nu}.
\end{equation*}
The three-index Levi-Civita symbol appearing in these equations is the usual rotation-
invariant symbol in three dimensions, with indices raised by the Euclidean 3-metric.

The Weyl equation implies that $p^2=0$ and that the solution with positive energy
has negative helicity (is left-handed) while the negative-energy solution has positive
helicity. The positive- and negative-helicity solutions are constrained two-component
spinors $\nu_\pm(p)$, with only one free component.

% BANKS-SOURCE: pdf=56 print=46 kind=display id=5.0b
In this $D=4$ setting, the general solution of the Weyl equation, with positive-frequency modes written as
$\ee^{+\ii p\cdot x}$, where $p\cdot x=-|\mathbf p|t+\mathbf p\cdot\mathbf x$, is
\begin{equation*}
  \nu(x)=\int\frac{\dd^3\mathbf p}{\sqrt{(2\pi)^3|\mathbf p|}}
  \left[a(\mathbf p)\ee^{+\ii p\cdot x}\nu_-(p)+b^\dagger(\mathbf p)\ee^{-\ii p\cdot x}\nu_+(p)\right].
\end{equation*}
Note that it is inconsistent to insist that $a$ and $b$ be the same operator, because the field
$\nu$ is complex. Quantization of $a$ and $b$ as either bosons or fermions is consistent with
the requirement that the Heisenberg equations derived from the Hamiltonian\footnote[6]{We are here anticipating Noether’s theorem from Chapter 6. The point is that time-translation symmetry allows us to construct the Hamiltonian from the Lagrangian without using the canonical formalism. Canonical commutators or anti-commutators can be \emph{derived} by insisting that Heisenberg’s equations be consistent with the Euler--Lagrange equations.} be the

% BANKS-SOURCE: pdf=57 print=47 kind=prose
Weyl equation. In $D=4$, the covariant PDE Green function of the Weyl operator, normalized by
$(-\ii\bar\sigma^\mu\partial_\mu)G_W(x)=\delta^4(x)$, can be constructed
without thinking about how the fields are quantized. The corresponding time-ordered
correlator is $\ii G_W$, and
\begin{equation*}
  G_W(x)=-\int\frac{\dd^4p}{(2\pi)^4}\ee^{+\ii p\cdot x}
  \frac{\sigma^\mu p_\mu}{p^2-\ii\epsilon}.
\end{equation*}
It is odd under interchange of spin indices combined with $x\to-x$ (one should write
things out in terms of the four real components of the Majorana spinor to see the
anti-symmetry). Thus, it is compatible only with Fermi statistics. Correspondingly (see
Problem 5.1) the Weyl field cannot have local commutation relations for either choice
of statistics, but has local anti-commutation relations for Fermi statistics. Thus, spin-$\frac12$
particles \emph{must} be fermions. The proof of the spin-statistics theorem follows the same
pattern for all spins \cite{banks-ref-6,banks-ref-7,banks-ref-8}, by analyzing the symmetry properties of the Feynman Green
function.

It is worth noting that we cannot allow fields with local anti-commutation relations
in the Lagrangian density, which must commute with itself at space-like separation,
but that even functions of them are allowed. Indeed, Lorentz invariance also requires
that we have only even functions of half-integer spin fields in the Lagrangian, so the
spin-statistics connection replaces these two constraints by one. It can be stated in
symmetry terms. Let $(-1)^F$ be the operator which is $-1$ on all states created from the
vacuum by products of fields containing an odd number of fermion operators and $+1$
on the other states. Then the spin-statistics connection is the equation
\begin{equation*}
  (-1)^F=\ee^{2\pi\ii\hat{\mathbf n}\cdot\mathbf J}.
\end{equation*}
($\hat{\mathbf n}$ is any unit 3-vector and $\mathbf J$ is the generator of angular momentum in any Lorentz
frame). The operator $(-1)^F$ commutes with all observables (operators that can be
viewed as infinitesimal changes of the Lagrangian). It should be emphasized that the
proof of the spin-statistics theorem follows from assuming that the Hilbert space of
the theory consists only of positive-norm particle states. Later on we will encounter
fermionic scalar fields, called Faddeev--Popov ghosts, which live in a ``Hilbert'' space of
indefinite norm. A gauge equivalence principle will forbid them from being produced in
the scattering of physical particles, but they are extremely useful at intermediate stages
of calculation.

% BANKS-SOURCE: pdf=57 print=47 kind=subsection id=5.0.1
\subsubsection{Chiral symmetry}

% BANKS-SOURCE: pdf=57 print=47 kind=prose
The Weyl Lagrangian has a symmetry under $\nu\to\ee^{\ii\alpha}\nu$, which gives rise via Noether’s
theorem (see Chapter 7) to a conservation law. At the level of free particles this con-
served quantity is simply helicity. This symmetry is called a chiral symmetry because
it acts differently on particles of different helicities. It cannot remain conserved for a
massive particle. Indeed, if we try to add a non-derivative term to the Weyl Lagrangian,
in order to generate a mass, then it must have the form
\begin{equation*}
  \mathcal{L}_m=m\epsilon^{\alpha\beta}\nu_\alpha\nu_\beta+\mathrm{h.c.},
\end{equation*}

% BANKS-SOURCE: pdf=58 print=48 kind=prose
where $m$ is a complex number. Note that this term would be identically zero if the fields
were $c$-numbers. Instead, in the next section we will show that they are anti-commuting
(Grassmann) numbers. In the presence of this term, the Weyl equation becomes
\begin{equation*}
  -\ii\bar\sigma^\mu\partial_\mu\nu+m^*\nu^*=0.
\end{equation*}
By multiplying this equation by $\sigma^\mu\partial_\mu$ and using the complex-conjugate equation, we
obtain
\begin{equation*}
  (\partial^2-mm^*)\nu=0,
\end{equation*}
indicating that the field creates and annihilates massive particles, with mass $|m|$.

The appearance of both $\nu$ and $\nu^*$ in the equation suggests that things will look
more elegant if we introduce the four-component Majorana field $\psi$. In terms of $\psi$, the
massive field equation takes the form
\begin{equation*}
  \left[-\ii\gamma^\mu\partial_\mu-mP_- -m^*P_+\right]\psi=0.
\end{equation*}
Here we have introduced the Dirac matrices,
\begin{equation*}
  \gamma^\mu\equiv\begin{pmatrix}0&\sigma^\mu\\ \bar\sigma^\mu&0\end{pmatrix}.
\end{equation*}
$P_\pm$ are the projectors on spinors of fixed chirality: $P_-$ projects on the first two
components of a Dirac spinor. In $D=4$, note that
\begin{equation*}
  \gamma_5=P_+-P_-=\begin{pmatrix}-1&0\\0&1\end{pmatrix}
\end{equation*}
anti-commutes with all the $\gamma^\mu$ and that $P_\pm=\frac12(1\pm\gamma_5)$. We define $\slashed A\equiv A_\mu\gamma^\mu$ for any
4-vector $A_\mu$. In the problems and Appendix C, you will find many properties of the
Dirac matrices. They mostly follow from the anti-commutation relation
\begin{equation*}
  [\gamma^\mu,\gamma^\nu]_+=-2\eta^{\mu\nu},
\end{equation*}
which the kind and diligent reader will easily verify.

We can do a chiral transformation on $\nu$ or $\psi$ to make the mass of a single free fermion
real. The Lagrangian for the Majorana field with real mass is

% BANKS-SOURCE: pdf=58 print=48 kind=equation id=5.2
\begin{equation}
  \mathcal{L}=\bar\psi(-\ii\slashed\partial-m)\psi,
  \label{eq:5.2}
\end{equation}
where\footnote[7]{The reader should prove that the $\gamma^0$ is necessary for the Lorentz invariance of $\mathcal{L}$.} $\bar\psi=\psi^\dagger\gamma^0$. This is called the Dirac Lagrangian. The general solution of its
equations of motion does not satisfy the Majorana condition, but can be decomposed
into a pair of fields, which do. This is analogous to the breakup of a complex scalar
field into real and imaginary parts. In fact (Problem 5.2) we can change basis in such a
way that the Dirac matrices are all imaginary. In this basis the Dirac equation has real
solutions and the Majorana condition is just $\psi^*=\psi$.

\BanksImplicitHook{I-CH05-002}
